% RESUMO - PT
\begin{resumo}
  \vspace{-15pt}

  A inspeção de anomalias superficiais em chapas de rochas ornamentais é feita visualmente por operadores, e o critério que separa uma feição natural de um defeito comercial não está escrito em lugar nenhum: ele varia com a rocha, o cliente, o lote e o inspetor. Este trabalho parte da premissa de que esse critério é arbitrário por natureza e propõe torná-lo explícito, parametrizado e auditável, em vez de tentar eliminá-lo. Propõe-se o ARIA (Análise e Reconhecimento Inteligente de Anomalias), um \textit{pipeline} hierárquico de visão computacional em que um classificador identifica a litologia da chapa, um modelo de fundação guiado por texto atua como Professor, gerando polígonos de anomalia sem anotação humana, e uma rede de segmentação compacta, a ser treinada sobre esses polígonos, atua como Aluno na linha de produção. Esta versão implementa e avalia os dois primeiros estágios; o treinamento do Aluno é trabalho futuro. Em cada litologia, o critério vira uma lista de termos que descrevem defeitos, as sondas, e, para cada sonda, um nível mínimo de confiança, o limiar, calibrados em quatro imagens; o limiar sai de uma regra automática, a mesma para todas as litologias. O experimento compara quatro configurações do Professor, e a hipótese opõe essa regra aplicada a cada litologia separadamente e a todas juntas. A referência é uma auditoria cega, em que uma avaliadora julga cada marcação sem saber de que configuração ela veio, com o critério de decisão fixado antes dos julgamentos. Nas onze litologias com mais imagens, o autor ajustou o limiar da regra em todos os 45 pares de litologia e sonda, sempre para cima e, em mediana, 3,3 vezes; em sete delas, subiu-o até que a superfície sem defeito da chapa comum deixasse de ser marcada. O classificador de litologias, reproduzido de trabalho anterior, atingiu 96,72\% de acurácia na partição de teste original e 99,92\% com partições sorteadas por imagem.
  % TODO: acrescentar uma frase de resultado do experimento quando a auditoria for julgada.

  Palavras-chave: \palavraschaveemlinha
\end{resumo}


% RESUMO - EN
\begin{resumo}[Abstract]
  \vspace{-15pt}

  \begin{otherlanguage*}{english}
  Surface anomaly inspection of polished dimension stone slabs is performed visually by human operators, and the criterion separating a natural feature from a commercial defect is nowhere written down: it varies with the stone, the customer, the batch and the inspector. This work takes that criterion to be arbitrary by nature and proposes to make it explicit, parameterised and auditable, rather than trying to eliminate it. We propose ARIA, a hierarchical computer vision pipeline in which a classifier identifies the slab lithology, a text-prompted foundation model acts as a Teacher, producing anomaly polygons without human annotation, and a compact segmentation network, to be trained on those polygons, acts as a Student on the production line. This version implements and evaluates the first two stages; training the Student is left as future work. For each lithology, the criterion becomes a list of terms describing defects, the probes, and, for each probe, a minimum confidence level, the threshold, calibrated on four images; the threshold comes from an automatic rule, the same for every lithology. The experiment compares four Teacher configurations, and the hypothesis opposes that rule applied to each lithology separately and to all of them pooled. The reference is a blind audit, in which an evaluator judges each marking without knowing which configuration produced it, with the decision criterion fixed before any judgement. Across the eleven lithologies with the most images, the author adjusted the rule's threshold for all 45 lithology--probe pairs, always upwards and by a median factor of 3.3; in seven of them, it was raised until the defect-free surface of the typical slab was no longer marked. The lithology classifier, reproduced from previous work, reached 96.72\% accuracy on the original test split and 99.92\% with per-image random splits.
  % TODO: mirror the experiment result sentence added to the Portuguese abstract.

  Keywords: \inlinekeywords
\end{otherlanguage*}
\end{resumo}
